\documentclass{report}
\usepackage{amsmath,amssymb}
\setlength{\parindent}{0mm}
\setlength{\parskip}{1em}
\begin{document}
\begin{center}
\rule{6in}{1pt} \
{\large Todd Munson \\
{\bf The TAO Augmented Lagrangian Method for PDE-Constrained Optimization}}

Math and Computer Science Division \\ Argonne National Laboratory \\ 9700 S Cass Ave \\ Argonne \\ IL 60439
\\
{\tt tmunson@mcs.anl.gov}\\
Evan Gawlik\\
Jason Sarich\\
Stefan Wild\end{center}

Parameter estimation, inverse problems, and optimal control problems
frequently appear in physical applications and can be cast
as constrained optimization problems with partial differential equation
constraints. In particular, we are interested in solving
discretized optimization problems of the form
\[
\begin{array}{ll}
\displaystyle \min_{u,v} & f(u,v) \\
\mbox{subject to} & g(u,v) = 0,
\end{array}
\]
where the state variable $u$ is the solution to the discretized partial
differential equation defined by $g$ and parametrized by the design
variable $v$, and $f$ is the discretized objective function, whose
form depends on the application. Examples include estimating the
porosity and permeability parameters from observables such as the
piezometric head at monitoring wells for reactive flows,
to applications in optical tomography and electromagnetic imaging.

In this talk, we will discuss the linearly-constrained augmented
Lagrangian method available in the Toolkit for Advanced Optimization (TAO)
for solving PDE-constrained optimization problems. This method computes
the direction in two phases, a Newton direction to reduce the constraint violation
and a reduced-space direction to improve the augmented Lagrangian merit
function. The reduced-space direction is computed from a limited-memory
quasi-Newton approximation to the reduced Hessian matrix. This method
requires a minimal amount of information from the user to solve the
optimization problem. The computational time for the method is dominated
by computing the necessary directions using iterative methods and
preconditioners supplied by the user for solving the linearized forward
and adjoint problems. The parallel performance is strongly influenced
by the parallel performance of these solves.

To evaluate the performance of our linearly-constrained augmented Lagrangian
method, we use the Haber-Hansen test problems, which include elliptic,
parabolic, and hyperbolic partial differential equations, and their
corresponding discretizations. We implemented parallel versions of these
applications in
the Toolkit for Advanced Optimization to investigate the scalability
of the method on high-performance architectures. Results from these
studies will be presented.

Our algorithms and test problems are available in the latest release
of the Toolkit for Advanced Optimization, which also includes methods
for derivative-free optimization, unconstrained and bound-constrained
optimization problems using derivative information, and complementarity
problems and variational inequalities on high-performance architectures.


\end{document}
